\chapter{O Sacramento da Unção dos Enfermos}

``Pela santa Unção dos Enfermos e pela oração dos presbíteros, toda a Igreja encomenda os doentes ao Senhor, sofredor e glorificado, para que os alivie e os salve: mais ainda, exorta-os a que, associando-se livremente à paixão e morte de Cristo, concorram para o bem do povo de Deus.'' Com estas palavras do Concílio Vaticano~II, o Catecismo da Igreja Católica abre o tratado sobre o sacramento da Unção dos Enfermos, nos parágrafos 1499 a 1532. Não se trata de um rito de passagem reservado ao último momento da vida, nem de uma formalidade religiosa diante da morte inevitável; trata-se de um sacramento de cura pelo qual Cristo, médico das almas e dos corpos, continua a tocar os enfermos com sua graça, unindo-os à sua própria Paixão redentora e antecipando, na fé, a vitória definitiva sobre o pecado e a morte.

\section{Os fundamentos na economia da salvação (1499--1513)}

A doença e o sofrimento estiveram sempre entre os problemas mais graves que afligem a vida humana. Na doença, o homem experimenta sua incapacidade, seus limites, sua finitude. Qualquer enfermidade pode fazer entrever a morte. A doença pode levar à angústia, ao fechamento em si mesmo e, por vezes, ao desespero e à revolta contra Deus. Mas também pode tornar uma pessoa mais amadurecida, ajudá-la a discernir o que não é essencial e a voltar-se para o essencial. Muitas vezes a doença leva à busca de Deus e a um regresso a Ele. A experiência humana do sofrimento não é, portanto, unívoca: pode ser destruidora ou purificadora, conforme seja ou não vivida na fé.

O homem do Antigo Testamento vive a doença diante de Deus. É diante d'Ele que desabafa seu lamento, e é d'Ele, Senhor da vida e da morte, que implora a cura. A doença torna-se caminho de conversão, e o perdão de Deus dá início à cura. Israel faz a experiência de que a doença está, de modo misterioso, ligada ao pecado e ao mal, e que a fidelidade a Deus restitui a vida. O profeta entrevê que o sofrimento pode ter também um sentido redentor pelos pecados dos outros. Isaías anuncia que Deus fará vir para Sião um tempo em que perdoará todas as faltas e curará todas as doenças (cf.~Is 33, 24).

A compaixão de Cristo para com os doentes e suas numerosas curas de enfermos de toda espécie são um sinal claro de que ``Deus visitou o seu povo'' (Lc 7, 16) e de que o Reino de Deus está próximo. Jesus tem poder não somente para curar, mas também para perdoar os pecados; veio curar o homem em sua totalidade, alma e corpo: é o médico de que os doentes precisam. Sua compaixão com todos os que sofrem vai ao ponto de identificar-Se com eles: ``Estive doente e visitastes-Me'' (Mt 25, 36). Jesus pede muitas vezes aos doentes que acreditem, e serve-se de sinais --- saliva e imposição das mãos, lodo e lavagem --- para curar; por seu lado, os doentes procuram tocá-L'O: ``porque saía d'Ele uma força que a todos curava'' (Lc 6, 19). Por isso, nos sacramentos, Cristo continua a ``tocar-nos'' para nos curar.

Comovido pelo sofrimento humano, Cristo não somente Se deixou tocar pelos doentes, mas fez suas as misérias deles: ``Tomou sobre Si as nossas enfermidades e carregou com as nossas doenças'' (Mt 8, 17). Ele não curou todos os doentes; as curas que realizava eram sinais da vinda do Reino. Na cruz, Cristo tomou sobre Si todo o peso do mal e tirou ``o pecado do mundo'' (Jo 1, 29), do qual a doença não é mais que consequência. Pela sua paixão e morte na cruz, Cristo deu novo sentido ao sofrimento: desde então, ele pode configurar-nos com Ele e unir-nos à sua paixão redentora. Seguindo a Cristo, os discípulos adquirem uma nova visão da doença e dos doentes: ``partiram e pregaram que era preciso cada um arrepender-se. Expulsavam muitos demônios, ungiam com óleo numerosos doentes, e curavam-nos'' (Mc 6, 12-13). A Igreja dos Apóstolos conhece um rito próprio em favor dos enfermos, atestado por São Tiago: ``Alguém de vós está doente? Chame os presbíteros da Igreja para que orem sobre ele, ungindo-o com óleo em nome do Senhor. A oração da fé salvará o doente e o Senhor o aliviará; e, se tiver cometido pecados, ser-lhe-ão perdoados'' (Tg 5, 14-15). A Tradição reconheceu neste rito um dos sete sacramentos da Igreja.

\section{Quem recebe, quem administra e como se celebra (1514--1519)}

A Igreja crê e confessa que, entre os sete sacramentos, há um especialmente destinado a reconfortar os que se encontram sob a provação da doença grave: a Unção dos Enfermos. Este sacramento não é ``só dos que estão prestes a morrer''. O tempo oportuno para recebê-lo chegou certamente quando o fiel começa, por doença ou por velhice, a estar em perigo de morte. Se um doente que recebeu a Unção recupera a saúde, pode, em caso de nova enfermidade grave, receber outra vez este sacramento. No decurso da mesma doença, pode ser repetido se o mal se agrava. É conveniente receber a Unção dos Enfermos antes de uma operação cirúrgica importante, e o mesmo se diga das pessoas de idade, cuja fragilidade se acentua.

Somente os sacerdotes --- bispos e presbíteros --- são ministros da Unção dos Enfermos. É dever dos pastores instruir os fiéis acerca dos benefícios deste sacramento e exortá-los a que, com espírito de fé, o peçam a tempo. Os fiéis, por sua vez, devem animar os enfermos a chamarem o sacerdote para receberem este sacramento, e os doentes devem preparar-se para recebê-lo com boas disposições, com a ajuda do seu pastor e de toda a comunidade eclesial, convidada a rodear de modo especial os enfermos com suas orações e atenções fraternas. O isolamento do doente é, em si mesmo, contrário ao espírito deste sacramento, que é eclesial por natureza.

Como todos os sacramentos, a Unção dos Enfermos é uma celebração litúrgica e comunitária, quer tenha lugar no seio da família, quer no hospital ou na Igreja, para um só doente ou para um grupo deles. É muito conveniente que seja celebrada durante a Eucaristia. Se as circunstâncias a tal convidarem, a celebração pode ser precedida pelo sacramento da Penitência e seguida pelo da Eucaristia. Enquanto sacramento da Páscoa de Cristo, a Eucaristia deveria ser sempre o último sacramento da peregrinação terrestre, o ``Viático'' da ``passagem'' para a vida eterna. A Palavra e o sacramento formam um todo inseparável nesta celebração: a liturgia da Palavra, precedida de um ato penitencial, abre a celebração, e as palavras de Cristo e o testemunho dos Apóstolos despertam a fé do doente e da comunidade.

Os atos litúrgicos essenciais da celebração são: os presbíteros da Igreja impõem em silêncio as mãos sobre os enfermos; rezam por eles na fé da Igreja --- esta é a epiclese própria deste sacramento ---; então conferem a unção com óleo benzido, se possível, pelo bispo. O óleo é o sinal sacramental desta graça: ungir com óleo, na tradição bíblica e na prática médica antiga, era ao mesmo tempo curar e consagrar. Assim, estes atos litúrgicos indicam a graça que este sacramento confere aos doentes: não somente alívio físico, mas a força do Espírito para viver cristãmente a experiência da fragilidade humana.

\section{Os efeitos da celebração (1520--1523)}

A primeira graça deste sacramento é uma graça de reconforto, de paz e de coragem para vencer as dificuldades próprias do estado de doença grave ou da fragilidade da velhice. Esta graça é um dom do Espírito Santo que renova a confiança e a fé em Deus, e dá força contra as tentações do Maligno --- especialmente a tentação do desânimo e da angústia da morte. Esta assistência do Senhor pela força do seu Espírito visa levar o doente à cura da alma, mas também à do corpo, ``se tal for a vontade de Deus'' (cf.~Tg 5, 15). Além disso, ``se ele cometeu pecados, ser-lhe-ão perdoados'' (Tg 5, 15): o sacramento da Unção dos Enfermos possui, portanto, uma dimensão de perdão que o aproxima do sacramento da Penitência e completa sua ação.

Pela graça deste sacramento, o enfermo recebe a força e o dom de unir-se mais intimamente à Paixão de Cristo: ele é, de certo modo, consagrado para produzir frutos pela configuração com a paixão redentora do Salvador. O sofrimento, sequela do pecado original, recebe um sentido novo: transforma-se em participação na obra salvífica de Jesus. Esta transformação do significado do sofrimento não é uma ideologia consoladora; é uma realidade teológica fundada no mistério pascal: o Cristo que sofreu e ressuscitou é o mesmo que, pelo sacramento, une a si o enfermo, fazendo do seu sofrimento uma oferenda junto à sua própria oblação ao Pai.

Os doentes que recebem este sacramento, ``associando-se livremente à paixão e morte de Cristo, concorrem para o bem do povo de Deus''. Ao celebrar este sacramento, a Igreja, na comunhão dos santos, intercede pelo bem do doente. E o doente, por seu lado, pela graça deste sacramento, contribui para a santificação da Igreja e para o bem de todos os homens, pelos quais a Igreja sofre e se oferece, por Cristo, a Deus Pai. O sofrimento do cristão doente, unido ao de Cristo, possui, portanto, um valor eclesial e até universal: não é um fato privado, mas um ato de amor e de oferta que edifica o Corpo de Cristo.

Os efeitos do sacramento são, portanto, ao mesmo tempo pessoais e eclesiais: purificam o enfermo de seus pecados remanescentes, fortalecem sua alma, animam sua esperança, confortam seu corpo quando isto convém à salvação espiritual, e o inserem mais profundamente na solidariedade do Corpo Místico. A Igreja reconhece que não é chamada a suprimir o sofrimento humano à força de um milagre de cura, mas a transformá-lo pela graça, ajudando o enfermo a viver sua condição como uma Páscoa pessoal --- passagem, com Cristo, da morte para a vida.

\section{O Viático, último sacramento do cristão (1524--1525)}

Àqueles que vão deixar esta vida, a Igreja oferece, além da Unção dos Enfermos, a Eucaristia como Viático. Recebida neste momento de passagem para o Pai, a comunhão do Corpo e Sangue de Cristo tem um significado e uma importância particulares. É semente de vida eterna e força de ressurreição, segundo as palavras do Senhor: ``Quem come a minha Carne e bebe o meu Sangue tem a vida eterna; e Eu ressuscitá-lo-ei no último dia'' (Jo 6, 54). Sacramento de Cristo morto e ressuscitado, a Eucaristia é aqui o sacramento da passagem da morte para a vida, deste mundo para o Pai. A palavra ``Viático'' significa literalmente ``provisão de viagem'': é o alimento para a última e mais decisiva das jornadas.

Assim como os sacramentos do Batismo, da Confirmação e da Eucaristia constituem uma unidade chamada ``sacramentos da iniciação cristã'', também se pode dizer que a Penitência, a Santa Unção e a Eucaristia como Viático constituem, quando a vida do cristão chega ao seu termo, ``os sacramentos que preparam a entrada na Pátria'' ou os sacramentos com que termina a peregrinação. Há uma simetria profunda entre o começo e o fim da vida cristã: assim como a vida nova em Cristo foi inaugurada pelos sacramentos da iniciação, assim é acompanhada e consumada pelos sacramentos de cura e de partida. A Unção dos Enfermos completa a conformação com a morte e ressurreição de Cristo que o Batismo havia começado e a Confirmação fortalecido; e a Eucaristia como Viático é a porta pela qual o cristão entra, em Cristo, na vida sem fim.

A morte cristã não é um evento solitário e sem sentido, mas um ato eclesial rodeado pela solicitude da Igreja inteira, que sustenta o moribundo com seus sacramentos, suas orações e sua presença fraterna. A Igreja não abandona seus filhos diante do mistério da morte; ela os acompanha com os meios que Cristo lhe confiou, assegurando-lhes que ``nem a morte, nem a vida [...] nem nenhuma outra criatura poderá nos separar do amor de Deus, em Cristo Jesus, Senhor nosso'' (Rm 8, 38-39). O Viático é, assim, a expressão mais concreta e mais terna desta fidelidade da Igreja, que permanece ao lado do moribundo até o limiar da eternidade.

\section{Em síntese (1526--1532)}

``Alguém de vós está doente? Chame os presbíteros da Igreja, para que orem sobre ele, ungindo-o com óleo em nome do Senhor. A oração da fé salvará o doente e o Senhor o aliviará. E, se tiver cometido pecados, ser-lhe-ão perdoados'' (Tg 5, 14-15). Com estas palavras do apóstolo São Tiago, o Catecismo resume o fundamento escriturístico e o sentido essencial do sacramento da Unção dos Enfermos. A sua finalidade é conferir uma graça especial ao cristão que enfrenta as dificuldades inerentes ao estado de doença grave ou de velhice. O tempo oportuno para recebê-la chegou certamente quando o fiel começa a encontrar-se em perigo de morte, por doença ou por velhice. Todas as vezes que um cristão cai gravemente enfermo, pode receber a Santa Unção; e também quando, mesmo depois de a ter recebido, a doença se agrava.

Somente os sacerdotes --- presbíteros e bispos --- podem ministrar o sacramento da Unção dos Enfermos; para isso, empregarão óleo benzido pelo bispo ou, em caso de necessidade, pelo próprio presbítero celebrante. O essencial da celebração deste sacramento consiste na unção na fronte e nas mãos do doente --- no rito romano --- ou sobre outras partes do corpo --- no Oriente ---, acompanhada da oração litúrgica do sacerdote celebrante que pede a graça especial deste sacramento. A forma sacramental proclama a misericórdia de Deus que vem ao encontro do enfermo: ``Por esta santa unção e pela sua infinita misericórdia, o Senhor venha em teu auxílio com a graça do Espírito Santo, para que, liberto dos teus pecados, Ele te salve e, na sua bondade, alivie os teus sofrimentos.''

A graça especial do sacramento da Unção dos Enfermos tem como efeitos: a união do doente à Paixão de Cristo, para seu bem e para o de toda a Igreja; o conforto, a paz e a coragem para suportar cristãmente os sofrimentos da doença ou da velhice; o perdão dos pecados, se o doente não pôde obtê-lo pelo sacramento da Penitência; o restabelecimento da saúde, se tal for conveniente para a salvação espiritual; e a preparação para a passagem para a vida eterna. Nenhum desses efeitos é mecânico ou automático: supõem a fé do enfermo e da comunidade que ora com ele, e realizam-se segundo os desígnios da providência divina, que sabe o que é melhor para cada um.

A Igreja cuida dos enfermos não por obrigação institucional, mas pelo amor de Cristo que, em cada doente, continua a revelar-Se e a convocar os cristãos ao serviço. O sacramento da Unção dos Enfermos é, por isso, ao mesmo tempo um sinal da predileção de Cristo pelos que sofrem e uma missão confiada à Igreja: ``Curai os enfermos!'' (Mt 10, 8). Mais do que um rito de consolação, ele é uma afirmação pascal --- a afirmação de que a última palavra não é a doença, nem a morte, mas a misericórdia de Deus que ressuscitou Jesus dos mortos e há de ressuscitar também aqueles que n'Ele acreditaram e que, pelo sacramento, se uniram à sua Paixão redentora.
